\documentclass[11pt,a4paper]{article}
\usepackage[T1]{fontenc}
\usepackage[margin=23mm]{geometry}
\usepackage{lmodern,mathrsfs,amsmath,amssymb,amsthm,mathtools,microtype,xurl}
\usepackage[hidelinks]{hyperref}
\newtheorem{theorem}{Theorem}
\title{Supported cohomology and counting-preserving\texorpdfstring{\\}{} inverses for ES divisor boxes}
\author{The Clankers}
\date{13 September 2026 --- research draft}
\begin{document}
\maketitle
\begin{abstract}
A finite three-term cochain identifies common residues without counting
internal cancellation as an arithmetic witness. Its original counting
norm has an explicit least-norm quotient. A different, fully marked
quotient recovers original divisor states from overlapping split boxes,
with inverse-fibre weights and actual relation primitives. These two
quotients need not have the same dimension. Iterated inverses agree when
the inherited metric is retained, and need not agree when it is reset.
The construction is a finite arithmetic instance of the existing SplitZero
quotient programme, not a new general cohomology or determinant theorem.
\end{abstract}

Let $p\equiv1\pmod4$ be prime, $p/4<a<p/2$, $R=4a-p$, and
$a=\prod q_i^{e_i}$. For a budget $b\ge0$ define
$B_b=\prod[-b_i,b_i]_{\mathbb Z}$ and
$v(\beta)=\prod q_i^{\beta_i}\pmod R$.
A channel label $c$ is M or E, with $t_M=-1$, $t_E=-p^{-1}$.
The divisor map $\beta\mapsto u=\prod q_i^{e_i+\beta_i}$ is a bijection
$B_e\to\operatorname{Div}(a^2)$. The two gates are $R\mid u+a$ and
$R\mid4u+1$. These are the inherited ES coordinates [4,5].

Choose $e=f+b$. Retain the two maps, separately at each channel,
\[
 X=B_f\xrightarrow{\ell}G,\quad Y=B_b\xrightarrow{r_c}G,
 \qquad\ell(\xi)=v(\xi),\quad r_c(\eta)=t_cv(\eta)^{-1},
 \quad G=(\mathbb Z/R\mathbb Z)^\times.                       \tag{1}
\]
For any finite maps $X\to Z\leftarrow Y$, let $B_X,B_Y$ be the free
star-difference modules in each nonempty fibre: choose reference $x_z$,
then use $e_x-e_{x_z}$ for every other $x$. Define
\[
 B_X\oplus B_Y\lhook\joinrel\longrightarrow k^X\oplus k^Y
 \xrightarrow{d}k^Z,
 \qquad d(A,B)_z=\sum_{X_z}A_x-\sum_{Y_z}B_y               \tag{2}
\]
in degrees $-1,0,1$ over any commutative ring $k$.

\begin{theorem}[Cohomology and original counting norm]
Put $I=\ell(X)\cap r(Y)$. The incoming map in (2) is injective,
$H^0(2)\simeq k^I$, and $H^1(2)\simeq k^{Z\setminus(\ell(X)\cup r(Y))}$.
The degree-zero map sends a cycle to its common fibre sums $c_z$; its
integral inverse is $c_z(e_{x_z},e_{y_z})$.
For the original counting Gram over $\mathbb R$ or $\mathbb C$, its
canonical representative assigns $c_z/m_z$ to each of the $m_z=|X_z|$
left occurrences and $c_z/n_z$ to each of the $n_z=|Y_z|$ right occurrences.
Its quotient squared norm is
\[
 \boxed{\sum_{z\in I}(m_z^{-1}+n_z^{-1})|c_z|^2.}             \tag{3}
\]
The source, relation and quotient determinants on that fibre are
$m_z+n_z$, $m_zn_z$, and $(m_z+n_z)/(m_zn_z)$.
\end{theorem}
\begin{proof}
A vector with zero fibre sum is uniquely
$\sum_{x\ne x_z}A_x(e_x-e_{x_z})$. In a cycle, only a common nonempty
fibre permits a nonzero sum. Subtracting the displayed reference lift
therefore gives exactly the incoming image and its explicit primitive.
The image of $d$ is the free coordinate module on the union of the two
images, proving $H^1$. The stated canonical representative has the same
sums and is orthogonal to every star difference; subtracting it proves
least norm and the Pythagorean identity. The determinant of a star Gram
is its fibre size. Applying the Schur identity to the cycle basis
(reference, star differences) gives the displayed determinants.
\end{proof}

The existing SplitZero reconstruction [1,2] retains a support mask and
distinguishes external absence $\tau$ from $(\lambda,0)$ in a present
fibre. Apply it to (2) at the masks of its two legs. Incoming relations
map to their \emph{supported} zero classes, not to absence. The source
quotient metric is the same actual-Gram construction as [3,\S1]. There is
no assertion that the finite counting norm equals the analytic theta
source norm.

A residue coarsening $\pi:Z\to Z'$ induces on $H^0$ the map
$c\mapsto(\sum_{\pi(z)=w}c_z)_w$. Its kernel has star-difference bases
inside fibres of $I\to\pi(I)$; its cokernel is freely indexed by
$(\pi\ell(X)\cap\pi r(Y))\setminus\pi(I)$. Every killed class has the
explicit coarser star primitive above. Thus local cohomology classes do
not assert one simultaneous original exponent choice.

\begin{theorem}[Original arithmetic masses and inverse fibres]
Let $\mathscr P_c=\{(\xi,\eta):\ell(\xi)=r_c(\eta)\}$ and
$\mathscr W_c=\{\beta\in B_e:v(\beta)=t_c\}$.
Addition $q:\mathscr P_c\to\mathscr W_c$ is onto, and its complete fibre is
\[
 \max(-f_i,\beta_i-b_i)\le\xi_i\le\min(f_i,\beta_i+b_i),
 \qquad\eta=\beta-\xi.                                   \tag{4}
\]
Let $N(\beta)$ be the product of these interval lengths.
For every original rational weight $w_\beta$,
\[
 \boxed{\sum_{\beta\in\mathscr W_c}w_\beta U^{e+\beta}
 =\sum_{(\xi,\eta)\in\mathscr P_c}
       \frac{w_{\xi+\eta}}{N(\xi+\eta)}U^{e+\xi+\eta}.}     \tag{5}
\]
The linear fibre-sum map $Q:k^{\mathscr P_c}\to k^{\mathscr W_c}$ has
exact star relations and an integral reference section. Over $\mathbb Q$
its counting-norm canonical section is $J(c)_x=c_{q(x)}/N(q(x))$.
The full inverse retains $QA$ and every nonreference coefficient of
$A-JQA$. Its quotient norm is $\sum|c_\beta|^2/N(\beta)$.
\end{theorem}
\begin{proof}
Each interval in (4) is nonempty since $|\beta_i|\le f_i+b_i$; its
bounds are exactly the two summand bounds. The residue product proves
surjectivity and (5). For the linear assertion, on each fibre sum-zero
vectors are exactly its star relations, and equal coefficients minimize
the counting norm. The relation primitive is the list of nonreference
coefficients. All inverse compositions follow by substitution.
\end{proof}

The original uniform norm on $k^{\mathscr W_c}$ is \emph{not} this native
quotient norm. Transporting it requires diagonal cost $N(\beta)$ on each
refined occurrence. This choice is forced among fibre-permutation-invariant
diagonal metrics: cost $h$ gives quotient cost $h/N(\beta)$.
Every change of metric is consequently explicit.

For finite surjections $P\to W\to Z$, put $m_w=|P_w|$ and
$n_z=\sum_{w\mapsto z}m_w$. With the inherited metric $1/m_w$, the second
canonical inverse gives $m_wc_z/n_z$ at $w$, and the first gives $c_z/n_z$
at every original occurrence. This is exactly the direct inverse. Resetting
the intermediate metric changes it. For three copies of $[-1,1]$, total
exponent zero has seven lifts. Direct coefficients are all $1/7$, squared norm
$1/7$; resetting yields four $1/6$ and three $1/9$, squared norm $4/27$. The
nonzero boundary discarded by resetting has squared norm
\[
 \boxed{4/27-1/7=1/189.}                                  \tag{6}
\]

\paragraph{The same object keeps failures and successful denominators.}
For any pair put $x=v(\xi)$, $w=v(\eta)$, $\rho=x-t_cw^{-1}$ and
$\beta=\xi+\eta$. The original gates obey
\[
 g_M=u+a\equiv aw\rho,\quad g_E=4u+1\equiv pw\rho\pmod R.
\]
Both factors are units, so $D=R/\gcd(R,g_c)=R/\gcd(R,\rho)$.
Formula (5) extends to \emph{all} pairs by adding $Z^{D-1}$ to each
monomial and using $\mathscr W_c=B_e$. All original prime-power defects
are retained. Multiplication by the same unit gives the explicit affine
coefficient-system conjugacy of [4], with inverse multiplication by its
modular inverse.

For a matched pair, set $d_0=\gcd(a,u)$,
$(h,r,s)=(d_0^2/u,u/d_0,a/d_0)$. On M let
$\lambda=(r+s)/R$ and take $(a,phs\lambda,phr\lambda)$.
On E let $\kappa=(pr+s)/R$ and take $(a,hs\kappa,phr\kappa)$.
The gate makes the quotient a positive integer; clearing reciprocals
proves $4/p$. The ordered inverse is $u=pa^2/(RY-pa)$ on M and
$u=a^2/(RY-pa)$ on E. Divisor valuations recover $\beta$, and (4)
recovers every split lift. The cohomology rank counts residues, not pairs
or original divisor states; those two projections of $\mathscr P_c$
remain separate.

\paragraph{Exact examples and nonclaims.}
At $p=2521,a=636,R=23$, split $(2,1,1)=(1,1,1)+(1,0,0)$.
There are five matched pairs, four channel/residue classes, and three
original divisor states. Four pairs have $N=2$, one has $N=1$, so the
original mass is $4/2+1=3$. The harmonic Gram is
$\operatorname{diag}(2,2,3/2,2)$, determinant $12$.
At the empty shell $p=241,a=64,R=15$, split $6=3+3$.
The full two-channel complex has dimensions $(12,28,16)$ and is exact.
The raw two-leg kernel has dimension twelve, all relations, while the
all-ones source norm remains $28$. A supported zero cohomology fibre is
not an absent source. Neither example is new prime coverage.

The general proofs above are self-contained finite algebra. The delivered
standard-library checker verifies actual differentials, Gram determinants,
coarsening primitives, the full marked examples, all small map pairs,
all small overlapping splits and a bounded arithmetic scan. It does not
prove universal ES or a zeta spectral identification. No historical
priority or Lean certification is claimed.

\begingroup\footnotesize
\begin{thebibliography}{5}
\bibitem{a} The Clankers (2026). \emph{SplitZero: reconstruction, internal homology, and balanced finite jets}, \S\S1--6. Zeta workbench, \texttt{formal/splitzero/DERIVED\_MATHEMATICS.md}, revision \texttt{7ea0a49945390eae14d3160a5730858899768b5f}.
\bibitem{b} The Clankers (2026). \emph{Arithmetic cohomology over the tau-base}, \S\S1--4. Same repository/revision, \texttt{workbenches/tau-base-cohomology/TAU\_BASE\_MODEL.md}.
\bibitem{c} The Clankers (2026). \emph{Canonical quotient volumes, trace control and an endpoint-budget criterion}, \S\S1,5. Same repository/revision, \texttt{workbenches/tau-toda-volume-formal/RESEARCH\_NOTE.md}.
\bibitem{d} The Clankers (2026). \emph{Split divisor boxes and exact Erd\H{o}s--Straus forcing}. ES workbench draft PR 5, revision \texttt{4f17e4aa63113a063eb027c8931cee5d5cce3981}.
\bibitem{e} Elsholtz, C., \& Tao, T. (2013). Counting the number of solutions to the Erd\H{o}s--Straus equation on unit fractions. \emph{J. Aust. Math. Soc., 94}(1), 50--105. \S2, Propositions 2.2 and 2.6.
\end{thebibliography}
\endgroup
\end{document}
